\chapter{Dati e preprocessing}
\label{chap:dati}

Il capitolo descrive i campi meteorologici su cui il detector è addestrato e valutato: le variabili della rianalisi ERA5, il dominio e la cadenza temporale, la preparazione dei dati e la loro suddivisione in anni di addestramento, di validazione e di test. Le trasformazioni che dipendono dal modello, come la normalizzazione stagionale, sono descritte nel Capitolo~\ref{chap:metodologia}; le osservazioni di stazione usate per la valutazione sono descritte nel Capitolo~\ref{chap:esperimenti}.

\section{Campi ERA5, allineamento e celle valide}
\label{sec:dati-era5}

I dati della rianalisi ERA5~\cite{hersbach2020era5} provengono da due prodotti distinti: variabili a livello singolo e variabili su livelli di pressione, di cui è usato soltanto il livello di 850 hPa. Le variabili sono quindici. Le dieci variabili a livello singolo sono temperatura a due metri, componenti orizzontali del vento a dieci metri, pressione al livello medio del mare, precipitazione totale, copertura nuvolosa totale, radiazione solare discendente alla superficie, umidità volumetrica del primo strato del suolo, velocità d'attrito e altezza dello strato limite. Le altre cinque sono le componenti orizzontali del vento, la velocità verticale, la temperatura e l'umidità specifica a 850 hPa. La Tabella~\ref{tab:variabili-era5} le elenca con le unità di ERA5, che non vengono convertite.

\begin{table}[htbp]
    \centering
    \footnotesize
    \caption{Le quindici variabili ERA5, con la sigla e l'unità del dataset. Precipitazione e radiazione sono accumuli di un'ora; le altre variabili sono istantanee.}
    \label{tab:variabili-era5}
    \vspace{0.6em}
    \begin{tabular}{@{}L{0.46}L{0.10}L{0.16}L{0.14}@{}}
        \toprule
        Variabile & Sigla & Livello & Unità \\
        \midrule
        Temperatura a due metri & t2m & Superficie & K \\
        Componente zonale del vento a dieci metri & u10 & Superficie & $\mathrm{m\,s^{-1}}$ \\
        Componente meridionale del vento a dieci metri & v10 & Superficie & $\mathrm{m\,s^{-1}}$ \\
        Pressione al livello medio del mare & msl & Superficie & Pa \\
        Precipitazione totale & tp & Superficie & m \\
        Copertura nuvolosa totale & tcc & Superficie & 0--1 \\
        Radiazione solare discendente alla superficie & ssrd & Superficie & $\mathrm{J\,m^{-2}}$ \\
        Umidità volumetrica del primo strato del suolo & swvl1 & Superficie & $\mathrm{m^{3}\,m^{-3}}$ \\
        Velocità d'attrito & zust & Superficie & $\mathrm{m\,s^{-1}}$ \\
        Altezza dello strato limite & blh & Superficie & m \\
        \addlinespace
        Componente zonale del vento & u & 850 hPa & $\mathrm{m\,s^{-1}}$ \\
        Componente meridionale del vento & v & 850 hPa & $\mathrm{m\,s^{-1}}$ \\
        Velocità verticale & w & 850 hPa & $\mathrm{Pa\,s^{-1}}$ \\
        Temperatura & t & 850 hPa & K \\
        Umidità specifica & q & 850 hPa & $\mathrm{kg\,kg^{-1}}$ \\
        \bottomrule
    \end{tabular}
\end{table}

I dati coprono gli anni dal 1980 al 2025. Sono selezionati gli istanti UTC 00, 03, 06, 09, 12, 15, 18 e 21, su una griglia di $0.5^\circ\times0.5^\circ$ nell'area compresa fra $20^\circ$ e $60^\circ$ di latitudine nord e fra $15^\circ$ di longitudine ovest e $50^\circ$ di longitudine est (Figura~\ref{fig:dominio-era5}). Il dominio comprende l'Europa, il Mediterraneo e parte dell'Africa settentrionale e del Medio Oriente. Il campionamento ogni tre ore è una selezione di tempi validi, non un'aggregazione temporale: precipitazione e radiazione conservano gli accumuli di un'ora associati ai tempi selezionati, senza somme triorarie o conversioni locali delle unità.

I file mensili vengono verificati rispetto a tempi, griglia, variabili e livello di pressione richiesti. La preparazione allinea i due prodotti sulle stesse coordinate e sugli stessi istanti, e ordina le latitudini da nord a sud e le longitudini da ovest a est. Il reticolo contiene $81\times131=10611$ celle.

\begin{figure}[htbp]
    \centering
    \includegraphics{dominio_era5}
    \caption{Dominio dei campi ERA5: $81\times131$ celle di $0.5^\circ$, fra $20^\circ$ e $60^\circ$ di latitudine nord e fra $15^\circ$ di longitudine ovest e $50^\circ$ di longitudine est. Il rettangolo indica il Piemonte, la regione dei dati del lavoro preliminare della Sezione~\ref{sec:rw-pvgis}. Linee di costa da Natural Earth.}
    \label{fig:dominio-era5}
\end{figure}
% Confini del dominio riletti il 21 ottobre 2026 nei dati preparati usati da
% tutte le run, data/era5_prepared/metadata.json (campo grid): 81 latitudini da
% 60 a 20, 131 longitudini da -15 a 50. Figura generata da
% figures/genera_mappa_era5.py; la linea di costa (Natural Earth, pubblico
% dominio, scala 1:50m) non è nel repository. Il rettangolo è l'area delle
% località PVGIS: se la figura viene rigenerata, si può togliere.

Sono conservate le celle per le quali tutte le variabili sono definite nel primo istante di addestramento, e si verifica che questo insieme rimanga identico nel resto del periodo; se cambia, la preparazione si interrompe. La selezione non utilizza il test e non viene applicata imputazione. Il criterio è la completezza delle variabili, non una classificazione terra--mare: nei dati usati nessuna cella è esclusa.

\section{Anni di addestramento, di validazione e di test}
\label{sec:dati-split}

I dati preparati hanno forma $[T,N,15]$, con $N=10611$ celle, e sono divisi in tre periodi consecutivi. Gli anni 1980--1999 sono usati per l'addestramento, gli anni 2000--2004 per la validazione e gli anni 2005--2025 per il test: rispettivamente 58440, 14616 e 61360 istanti. Le quantità stimate dai dati per l'ingresso e la loss del modello, cioè la normalizzazione stagionale, la tendenza rimossa, gli estremi del min--max e i pesi delle variabili del Capitolo~\ref{chap:metodologia}, sono calcolate sui soli anni di addestramento.

La validazione serve a confrontare le configurazioni e a fissare, prima di valutare il test, tutte le quantità che descrivono il comportamento abituale degli score. Queste quantità sono la climatologia degli score, gli estremi con cui le due componenti dello score combinato sono riportate sulla stessa scala, le scale e la climatologia dei rilevatori senza modello, le soglie e i limiti delle classi di fiducia (Capitolo~\ref{chap:esperimenti}). Sulla validazione è fatta anche la scelta dello score. Gli anni di test non entrano in alcuna scelta: ricevono senza modificarle le quantità stimate sugli anni di addestramento e su quelli di validazione. L'articolo di STGAN~\cite{deng2022stgan} separa i dati in training e test senza un insieme di validazione: la validazione è un'aggiunta di questo lavoro, necessaria perché in assenza di etichette la scelta della configurazione deve basarsi su una misura calcolata fuori dal test.

I primi istanti di ciascun periodo usano come contesto la coda del periodo precedente: le finestre di ingresso dei primi istanti di validazione contengono gli ultimi istanti di addestramento, e lo stesso vale fra validazione e test. Il contesto non rientra fra gli istanti valutati.
% Conteggi dal calendario, a otto istanti al giorno: 1980-1999 sono 7305 giorni
% (cinque bisestili), 2000-2004 sono 1827 giorni, 2005-2025 sono 7670 giorni.
% 61360 istanti di test e 10611 celle anche dalle dimensioni degli score di una
% run completa. Il numero di aggiornamenti per epoca dello sweep (58384) è
% 58440 meno i 56 istanti di contesto del trend.
% ATTENZIONE: i risultati provvisori del Capitolo 6 vengono ancora dalla run
% ..._wdaily_seasonal_trend_ep3 (suddivisione precedente). La run della tesi è
% good-sweep-2 (sweep 913ez3rt, W&B 5qlj2lwg), con la suddivisione qui descritta:
% i risultati vanno riscritti con i suoi numeri (outputs/good2v, analisi_good2).

\section{Caratteristiche dei campi}
\label{sec:dati-caratteristiche}

Tre proprietà dei campi motivano altrettante scelte del Capitolo~\ref{chap:metodologia}. Sono misurate sugli anni di addestramento, senza alcun modello.

\paragraph{Cicli diversi da luogo a luogo.} La temperatura a due metri e la copertura nuvolosa hanno un ciclo stagionale e uno giornaliero, con livello e ampiezza diversi da una cella all'altra (Figura~\ref{fig:era5-cicli}). La media mensile della temperatura va da 1 a 17~°C in una cella dell'Europa settentrionale, da 2 a 24~°C nella Pianura Padana e da 11 a 32~°C nel Sahara; la copertura nuvolosa media sta fra 0.60 e 0.78 nella prima cella e fra 0.04 e 0.29 nell'ultima. Lo stesso valore è quindi abituale in un luogo e in una stagione e insolito in altri. Per questo i valori in ingresso sono riferiti al normale della propria cella, del giorno e dell'ora (Sezione~\ref{sec:ingresso-normalizzazione-era5}), e gli score alla climatologia della propria cella (Sezione~\ref{sec:normalizzazione-climatologica-score}).

\begin{figure}[htbp]
    \centering
    \includegraphics{era5_cicli}
    \caption{Ciclo stagionale e giornaliero della temperatura a due metri e della copertura nuvolosa in tre celle: Europa settentrionale ($55^\circ$N, $10^\circ$E), Pianura Padana ($45^\circ$N, $9^\circ$E) e Sahara ($25^\circ$N, $10^\circ$E). A sinistra la media per mese, con una deviazione standard; a destra la media per ora UTC in gennaio e in luglio. Anni 1980--1999.}
    \label{fig:era5-cicli}
\end{figure}

\paragraph{Variazioni di grandezza diversa.} Con il min--max dell'articolo, che riporta ogni variabile fra il minimo e il massimo degli anni di addestramento, la copertura nuvolosa produce l'84.1\% della variazione quadratica media su 24 ore delle quindici variabili, e nessuna delle altre supera il 3.3\% (Figura~\ref{fig:era5-variazione}). La copertura nuvolosa passa spesso da sereno a coperto nel corso di un giorno e attraversa quasi tutto il proprio intervallo; la temperatura cambia di pochi gradi, in un intervallo che comprende tutte le celle e tutte le stagioni. Un errore di ricostruzione che somma le variabili con lo stesso peso sarebbe quindi in gran parte un errore sulla copertura nuvolosa. Dopo la normalizzazione stagionale e il min--max del detector le quote vanno dallo 0.9\% dell'umidità del suolo al 16.6\% dell'umidità specifica, e la copertura nuvolosa ha il 6.7\%, la quota di una variabile media: la normalizzazione stagionale rende confrontabili anche le variabili fra loro. Le differenze che restano sono corrette dai pesi della Sezione~\ref{sec:pesi-variabili-ricostruzione}.

\begin{figure}[htbp]
    \centering
    \includegraphics{era5_variazione_24h}
    \caption{Quota di ciascuna variabile nella variazione quadratica media su 24 ore, dopo il min--max dell'articolo e dopo la normalizzazione stagionale e il min--max del detector. Anni 1980--1999; la seconda quota è ricavata dai pesi della Sezione~\ref{sec:pesi-variabili-ricostruzione}. La linea tratteggiata indica la quota uguale per tutte le variabili, un quindicesimo.}
    \label{fig:era5-variazione}
\end{figure}

\paragraph{Legame spaziale a corto raggio.} Le anomalie di due celle, cioè gli scostamenti dalla media della propria cella, del mese e dell'ora, sono tanto meno correlate quanto più le celle sono lontane (Figura~\ref{fig:era5-correlazione}). La correlazione mediana scende sotto 0.5 oltre circa 1300~km per la pressione al livello del mare e 900~km per la temperatura a due metri, ma già oltre 200~km per la copertura nuvolosa e 100~km per la precipitazione. Le variabili legate alle perturbazioni dipendono quindi dalle celle vicine, entro poche celle di $0.5^\circ$: è il contesto che il detector rappresenta con la ConvGRU del generatore e con la patch del discriminatore (Sezione~\ref{sec:stgan}).

\begin{figure}[htbp]
    \centering
    \includegraphics{era5_correlazione_distanza}
    \caption{Correlazione delle anomalie di due celle in funzione della loro distanza: mediana e quartili per fasce di 100~km, su 300 celle scelte a caso e un istante ogni nove degli anni 1980--1999. Il primo punto, entro 100~km, si basa su 36 coppie di celle.}
    \label{fig:era5-correlazione}
\end{figure}
% Numeri da figure_dati.log e figures/dati/era5_dati.json (22 ottobre 2026,
% scripts/summarize_era5_data_figures.py --train-end-year 1999 del repository del
% codice); quote del detector da figures/dati/era5_pesi_good2.json (pesi salvati
% nei metadati di good-sweep-2, stimati su 720 istanti). Figure generate da
% figures/genera_figure_dati.py.
